\section{Discussion}

\noindent\textbf{Why the decoder is the bottleneck.}
SAM's ViT-H encoder, pre-trained on 1B masks, already captures texture gradients and edge patterns sufficient for COD. The decoder---trained on general segmentation---has never learned that \emph{subtle} texture disruptions signify boundaries. Fine-tuning teaches this mapping: subtle features $\rightarrow$ object boundaries. The 16.6-point mIoU gain from training only 0.6\% of parameters (4M of 636M) confirms this bottleneck hypothesis.

\noindent\textbf{Prompt robustness from mixed training.}
By alternating point and box prompts, the decoder learns \emph{what} the camouflaged object is rather than \emph{where the user clicked}. Point-only training overfits to spatial proximity; box-only training overfits to box geometry. Mixed training forces the decoder to understand object structure, explaining why improvement transfers to prompt types never seen during training.

\noindent\textbf{LLPM: targeted boundary enhancement.}
The LLPM's impact is concentrated on edge prompts (+2.6 mIoU) and dark/boundary-ambiguous images. The learned gate converged to $\alpha = 0.185$, a moderate perturbation that enhances boundaries without disrupting the global statistics the frozen encoder expects. This confirms that pre-encoder and decoder-side adaptation are complementary.

\noindent\textbf{Oracle prompts vs.\ automatic methods.}
Table~\ref{tab:comparison} contextualizes decoder-only adaptation within the broader COD landscape, but is not a head-to-head comparison: our method uses oracle GT-derived prompts while dedicated COD methods are fully automatic. Even with this advantage, we do not surpass encoder-adaptation methods on $S_\alpha$ or MAE, confirming that decoder-only fine-tuning is not a substitute for encoder modification. Rather, the comparison illustrates that decoder specialization provides a strong and efficient starting point. In practice, prompts come from interactive users or upstream detectors; the resulting performance will depend on prompt quality. A natural next step is pairing our decoder with an automatic prompt generator (e.g., a lightweight detection head) to enable fair comparison with fully automatic pipelines.

\noindent\textbf{Limitations.}
(1) Primarily evaluated on COD10K; cross-dataset experiments on CAMO show consistent generalization (+18.8 mIoU), with NC4K results pending. (2) Requires user prompts, unlike fully automatic COD methods. (3) Encoder-side adaptation (SAM-Adapter) achieves higher absolute performance. (4) LLPM training requires live encoder forward passes (7.5 hrs vs.\ 90 min for decoder-only).
